\section{Computational Results and Course Compliance}
\label{sec:results}

The first experimental objective is to verify that the complete course-style
workflow is available for Tracks. A manual instance named
\texttt{instanceTest.txt} is provided, and the solver can also be applied to
the existing small manual instances. For each instance, the result file records
at least the variables \texttt{solveTime} and \texttt{isOptimal}, matching the
format expected by the course example.

\begin{table}[h]
\centering
\caption{Representative run of the exact MILP solver on manual Tracks instances.}
\label{tab:manual-results}
\begin{tabular}{lrrrr}
\toprule
\textbf{Instance} & \textbf{Rows} & \textbf{Cols} & \textbf{Time (s)} & \textbf{Optimal} \\
\midrule
\texttt{instanceTest.txt} & 4 & 4 & 0.038 & yes \\
\texttt{small\_4x4.txt} & 4 & 4 & 0.040 & yes \\
\texttt{small\_5x5.txt} & 5 & 5 & 0.045 & yes \\
\texttt{unsat\_2x2.txt} & 2 & 2 & 0.036 & no \\
\bottomrule
\end{tabular}
\end{table}

The first three instances are feasible and are solved to optimality by the
MILP backend. The fourth instance is intentionally inconsistent and is reported
as infeasible. This is a useful check because it confirms that the solver does
not merely return a route whenever it is called; it also detects impossible
clue systems.

The generated course-style result files are stored in
\texttt{res/tracks/cbc}. A CSV summary is also produced for easier inspection,
and a LaTeX table can be generated automatically from the same result files.
The exact numerical times in Table~\ref{tab:manual-results} are machine
dependent, but the important point for this stage is the structure of the
workflow: every instance can be read, solved, validated, displayed, and
summarized in a reusable result format.

From the point of view of the project statement, the remaining experimental
work for Tracks is therefore not the construction of the model, but the
extension of the benchmark set. Larger generated instances and screenshot
instances can be solved with the same pipeline, then summarized with the same
result-table function. These experiments will allow the final report to discuss
how solving time changes with grid size and clue structure.
